\reading{MR-13}{Expectations, Risk Premium, Convexity, and the Shape of the Term Structure}{defer}{2}
% ==========================================================

\begin{mrkeybox}[title={Learning objectives in this reading}]
\small
\begin{enumerate}[label=\textbf{\alph*.},leftmargin=1.7em]
\item Explain the role of interest rate expectations in determining the shape of the term
      structure.
\item Apply a risk-neutral interest rate tree to assess the effect of volatility on the shape of
      the term structure.
\item Estimate the convexity effect using Jensen's inequality.
\item Identify the components into which the return on a bond can be decomposed, and calculate
      the expected return on a bond for a risk-averse investor.
\end{enumerate}
\end{mrkeybox}

% ---------------------------------------------------------
\lo{MR-13 a}{Explain the role of interest rate expectations in determining the shape of the term structure.}

\subsection{Three calculations you will need}

\begin{mrfmlbox}
\textbf{1. Spot rate given the price of a bond}
\[ \$0.85763 \;=\; \frac{\$1}{\bigl(1 + \hat{r}(2)\bigr)^{2}} \]
\vspace{-10pt}
\[ \text{PV} = -0.85763,\ \ \text{FV} = 1,\ \ \text{PMT} = 0,\ \ N = 2
   \ \Rightarrow\ I/Y = \mathbf{7.9816\%} \]
\vspace{-4pt}
\mrvk{$\hat{r}(2)$ = the implied 2-year spot rate. Solve the price equation for the rate, or use
the TVM keys.}{}

\medskip
\textbf{2. Spot rate given expected future spot rates}
\[ \frac{\$1}{(1.08)(1.10)} \;=\; \frac{\$1}{\bigl(1+\hat{r}(2)\bigr)^{2}}
   \qquad\longrightarrow\qquad
   \hat{r}(2) = \sqrt[2]{(1.08)(1.10)} - 1 = \sqrt{1.188} - 1 = \mathbf{8.995\%} \]
\vspace{-4pt}
\mrvk{Chain the one-year rates, then take the $n$th root. This is the pure expectations result
with \textit{no} volatility and \textit{no} risk premium.}{}

\medskip
\textbf{3. Price of a bond}
\[ \text{1-year zero} = \frac{\$1}{1.08} = \$0.92593
   \qquad\qquad
   \text{2-year zero} = \frac{\$1}{(1.08)(1.08)} = \$0.85734 \]
\end{mrfmlbox}

\subsection{How expectations shape the curve}

\begin{center}
\small
\begin{tabularx}{\linewidth}{@{}p{2.5cm} >{\raggedright\arraybackslash}X >{\raggedright\arraybackslash}X >{\raggedright\arraybackslash}X@{}}
\arrayrulecolor{FRMRule}\toprule
\rowcolor{FRMNavy} \hdr{} & \hdr{Flat yield curve} & \hdr{Upward-sloping} & \hdr{Downward-sloping}\\
\textbf{1-year rate} & 8\% & 8\% & 8\%\\
\rowcolor{FRMSoft}
\textbf{Expected 1-year rates} & 8\% for the next two years & 10\% and 12\% in one and two years
  & 6\% and 4\% in one and two years\\
\textbf{Spot rates} &
  1-year: 8\%\newline 2-year: 8\%\newline 3-year: 8\% &
  1-year: 8\%\newline 2-year: 8.995\%\newline 3-year: 9.988\% &
  1-year: 8\%\newline 2-year: 6.995\%\newline 3-year: 5.987\%\\
\rowcolor{FRMSoft}
\textbf{Result} & Yield curve is flat; investors are willing to lock in rates for two or three
  years at 8\%. & Upward-sloping, because expected future rates exceed the current rate. &
  Downward-sloping, because expected future rates are lower than the current rate.\\
\bottomrule
\end{tabularx}
\end{center}

\begin{center}
\begin{tikzpicture}
\begin{axis}[
  width=10.6cm, height=5.5cm,
  xmin=0.7, xmax=3.3, ymin=5.0, ymax=10.8,
  xlabel={\scriptsize Maturity (years)}, ylabel={\scriptsize Spot rate (\%)},
  xlabel style={font=\scriptsize}, ylabel style={font=\scriptsize},
  tick label style={font=\scriptsize},
  xtick={1,2,3}, ytick={5,6,7,8,9,10},
  legend style={font=\scriptsize\sffamily, draw=FRMRule, fill=white,
    at={(0.5,-0.30)}, anchor=north, legend columns=-1, column sep=10pt},
  axis line style={draw=black!55}, clip=false]
\addplot[FRMGreen, line width=1.2pt, mark=*, mark size=1.6pt]
  coordinates {(1,8)(2,8.995)(3,9.988)};
\addlegendentry{Rising expectations}
\addplot[black!65, line width=1.2pt, dashed, mark=*, mark size=1.6pt]
  coordinates {(1,8)(2,8)(3,8)};
\addlegendentry{Flat expectations}
\addplot[FRMRed, line width=1.2pt, dotted, mark=*, mark size=1.6pt]
  coordinates {(1,8)(2,6.995)(3,5.987)};
\addlegendentry{Falling expectations}
\end{axis}
\end{tikzpicture}
\end{center}
\figcap{All three curves start at the same 8\% one-year rate. Everything that happens after that
is driven purely by what the market expects future one-year rates to be. Under pure
expectations, nothing else is in the model.}

\begin{mrkeybox}[title={Pure expectations theory}]
\begin{enumerate}[label=\textbf{\alph*)}]
\item Under pure expectations theory, the term structure is a function of expected future spot
      rates, and the shape of the spot rate curve is determined \term{solely} by expected future
      spot rates.
\item Hence forecasts can be very useful in describing the shape and level of the term structure
      over short-term horizons, \term{but} probably only the \textit{level} of rates --- not the
      \textit{shape}.
\end{enumerate}
\end{mrkeybox}

% ---------------------------------------------------------
\lo{MR-13 b}{Apply a risk-neutral interest rate tree to assess the effect of volatility on the shape of the term structure.}

Now put \term{volatility} into interest rate expectations. Assuming risk neutrality, the decision
tree below gives the price of a 2-year zero-coupon bond with a face value of \$1 using the
expected rates.

\begin{center}
\begin{tikzpicture}[every node/.style={font=\scriptsize\sffamily},
  rt/.style={circle, draw=FRMNavy!55, fill=PaleNavy, line width=0.7pt, inner sep=2.5pt,
    minimum size=0.82cm},
  pv/.style={rectangle, rounded corners=2pt, draw=FRMGreen!55, fill=PaleGreen,
    line width=0.7pt, inner sep=3.5pt, minimum width=1.5cm, align=center},
  lf/.style={rectangle, rounded corners=2pt, draw=black!35, fill=FRMSoft,
    line width=0.6pt, inner sep=3pt, minimum width=0.75cm}]
% rate tree
\begin{scope}
  \node[rt] (r0) at (0,0) {8\%};
  \node[rt] (ru) at (1.9,1.15) {10\%};
  \node[rt] (rd) at (1.9,-1.15) {6\%};
  \node[lf] (a) at (3.9,2.0) {12\%};
  \node[lf] (b) at (3.9,0) {8\%};
  \node[lf] (c) at (3.9,-2.0) {4\%};
  \foreach \p/\q in {r0/ru, r0/rd, ru/a, ru/b, rd/b, rd/c}{
    \draw[FRMNavy!60, line width=0.7pt] (\p) -- node[sloped, above, text=black]{50\%} (\q);}
  \node[anchor=north, text=black!60] at (1.9,-2.8) {\textbf{Expected rates}};
\end{scope}
% price tree
\begin{scope}[shift={(8.3,0)}]
  \node[pv] (p0) at (0,0) {\$0.85763};
  \node[pv] (pu) at (2.2,1.15) {\$0.90909};
  \node[pv] (pd) at (2.2,-1.15) {\$0.94340};
  \node[lf] (x) at (4.3,2.0) {\$1};
  \node[lf] (y1) at (4.3,0.5) {\$1};
  \node[lf] (y2) at (4.3,-0.5) {\$1};
  \node[lf] (z) at (4.3,-2.0) {\$1};
  \foreach \p/\q in {p0/pu, p0/pd, pu/x, pu/y1, pd/y2, pd/z}{
    \draw[FRMGreen!60, line width=0.7pt] (\p) -- node[sloped, above, text=black]{50\%} (\q);}
  \node[anchor=north, text=black!60] at (2.2,-2.8) {\textbf{Resulting prices}};
\end{scope}
\draw[FRMGold, -{Stealth[length=5pt]}, line width=1pt] (5.1,0) -- (6.5,0);
\end{tikzpicture}
\end{center}
\figcap{The expected price in one year at the upper node is $\$1/1.10 = \$0.90909$; at the lower
node, $\$1/1.06 = \$0.94340$. Discounting the average back one more year at 8\% gives today's
price.}

\begin{mrexbox}
\small
\[ \bigl[0.5 \times (\$0.90909 / 1.08)\bigr] + \bigl[0.5 \times (\$0.94340 / 1.08)\bigr]
   \;=\; \mathbf{\$0.85763} \]
With the present value in hand, solve for the implied 2-year spot rate:
\[ \text{PV} = -0.85763,\ \text{FV} = 1,\ \text{PMT} = 0,\ N = 2
   \ \Rightarrow\ I/Y = \mathbf{7.9816\%} \]
\end{mrexbox}

\begin{mrkeybox}[title={Impact of volatility}]
When there is uncertainty regarding expected rates, the volatility of expected rates causes the
future spot rates to be \term{lower}. With the implied rate, the value of convexity for the
2-year zero-coupon bond is:
\[ 8\% - 7.9816\% \;=\; 0.0184\% \;=\; \mathbf{1.84 \text{ basis points}} \]
\end{mrkeybox}

% ---------------------------------------------------------
\lo{MR-13 c}{Estimate the convexity effect using Jensen's inequality.}

\begin{mrdefbox}[title={The convexity effect}]
\begin{enumerate}[label=\textbf{\alph*)}]
\item This effect considers the \term{curvature} of the relationship between bond price and
      interest rates.
\item The convexity effect arises from a special case of \term{Jensen's inequality}.
\item Jensen's inequality states that the expected value of a convex function of a random
      variable is \textit{greater than or equal to} the convex function of the expected value of
      the variable.
\end{enumerate}
\end{mrdefbox}

\begin{mrfmlbox}
\[ E\!\left[ \frac{1}{1+r} \right] \;>\; \frac{1}{E(1+r)} \]
\mrvk{$r$ = the uncertain future rate. The left side averages the \textit{discount factors}; the
right side discounts at the \textit{average rate}. Because $1/(1+r)$ is convex, averaging first
and inverting second gives a smaller number than inverting first and averaging second.}{}
\end{mrfmlbox}

\begin{center}
\begin{tikzpicture}
% ---- left panel: why the inequality holds ----
\begin{axis}[name=conc, width=6.5cm, height=5.2cm,
  xmin=0, xmax=0.52, ymin=0.62, ymax=1.03,
  xlabel={\scriptsize Rate $r$}, ylabel={\scriptsize $1/(1+r)$},
  xlabel style={font=\scriptsize}, ylabel style={font=\scriptsize},
  tick label style={font=\tiny},
  xtick={0,0.1,0.2,0.3,0.4,0.5},
  xticklabels={0,10\%,20\%,30\%,40\%,50\%},
  ytick={0.7,0.8,0.9,1.0},
  title={\scriptsize\sffamily\bfseries\color{FRMNavy}The shape, exaggerated},
  axis line style={draw=black!55}, clip=false]
\addplot[black!80, line width=1.1pt, smooth, domain=0:0.52, samples=120] {1/(1+x)};
\addplot[FRMRed, line width=1.0pt, dashed] coordinates {(0.05,0.952381)(0.45,0.689655)};
\node[circle, fill=black!55, inner sep=1.4pt] at (axis cs:0.05,0.952381) {};
\node[circle, fill=black!55, inner sep=1.4pt] at (axis cs:0.45,0.689655) {};
\node[circle, fill=FRMRed, inner sep=1.7pt] at (axis cs:0.25,0.821018) {};
\node[circle, fill=FRMGreen, inner sep=1.7pt] at (axis cs:0.25,0.8) {};
\draw[FRMGold, line width=0.8pt] (axis cs:0.25,0.8) -- (axis cs:0.25,0.821018);
\node[anchor=west, font=\tiny\sffamily, text=FRMGold, fill=white, inner sep=1.2pt]
  at (axis cs:0.27,0.845) {the gap};
\node[anchor=north, font=\tiny\sffamily, text=FRMRed, fill=white, inner sep=1.2pt]
  at (axis cs:0.13,0.90) {chord};
\node[anchor=south, font=\tiny\sffamily, text=black!70, fill=white, inner sep=1.2pt]
  at (axis cs:0.36,0.73) {curve};
\end{axis}
% ---- right panel: the gap at the example's own numbers ----
\begin{axis}[name=zoom, at={($(conc.east)+(1.5cm,0)$)}, anchor=west,
  width=6.5cm, height=5.2cm,
  xmin=0.058, xmax=0.102, ymin=0, ymax=38,
  xlabel={\scriptsize Rate $r$},
  ylabel={\scriptsize Gap, in units of $10^{-5}$},
  xlabel style={font=\scriptsize}, ylabel style={font=\scriptsize},
  tick label style={font=\tiny},
  xtick={0.06,0.07,0.08,0.09,0.10},
  xticklabels={6\%,7\%,8\%,9\%,10\%},
  ytick={0,10,20,30},
  title={\scriptsize\sffamily\bfseries\color{FRMNavy}The gap, at the example's numbers},
  axis line style={draw=black!55}, clip=false]
\addplot[FRMGold, line width=1.2pt, smooth] coordinates {
(0.060,0.0)(0.065,14.093)(0.070,24.046)(0.075,29.917)(0.080,31.764)
(0.085,29.642)(0.090,23.605)(0.095,13.706)(0.100,0.0)};
\draw[FRMRed, dashed, line width=0.7pt] (axis cs:0.08,0) -- (axis cs:0.08,31.764);
\node[circle, fill=FRMRed, inner sep=1.7pt] at (axis cs:0.08,31.764) {};
\node[anchor=north, font=\tiny\sffamily, text=FRMRed, fill=white, inner sep=1.2pt,
  align=center] at (axis cs:0.08,29.0)
  {$0.92624 - 0.92593$\\ $= 0.00031$};
\end{axis}
\end{tikzpicture}
\end{center}
\figcap{Left, on an exaggerated rate range so the shape is visible: the chord joining two rate
outcomes lies \textit{above} the curve at the midpoint, and that vertical gap is the convexity
effect. Right, the same gap computed at this example's own 6\% and 10\% outcomes. It peaks at
the midpoint, 8\%, at $0.00031$ --- which is exactly $\$0.92624 - \$0.92593$. Plotted on the
discount-factor axis the gap would be invisible, which is why it is plotted on its own.}

\begin{mrexbox}[title={Example --- applying Jensen's inequality}]
Assume that next year there is a 50\% probability that 1-year spot rates will be 10\% and a 50\%
probability that they will be 6\%. Demonstrate Jensen's inequality for a 2-year zero-coupon bond
with a face value of \$1.
\tcblower
\small
\textbf{Left-hand side} --- the expected price in one year using the 1-year spot rates of 10\%
and 6\%:
\[ E\!\left[\frac{\$1}{1+r}\right] = 0.5 \times \frac{\$1}{1.10} + 0.5 \times \frac{\$1}{1.06}
   = \$0.92624 \]

\textbf{Right-hand side} --- the expected price using an expected rate of 8\%:
\[ \frac{\$1}{0.5 \times 1.10 + 0.5 \times 1.06} = \frac{\$1}{1.08} = \$0.92593 \]

So the left-hand side is greater: $\$0.92624 > \$0.92593$.

\medskip
\textbf{Now discount one more year at 8\%.} The left-hand side becomes
$\$0.92624 / 1.08 = \mathbf{\$0.85763}$. The right-hand side is the 2-year zero discounted for
two years at 8\%, $\$1 / 1.08^2 = \mathbf{\$0.85734}$.

Again the left-hand side is greater: $\$0.85763 > \$0.85734$.

\medskip
Since the 2-year zero-coupon price is \textit{higher} than the price achieved through
straightforward discounting, its implied rate must be \term{lower} than 8\%.
\end{mrexbox}

\begin{mrkeybox}[title={Remember this relationship}]
\begin{center}
\begin{tikzpicture}[every node/.style={font=\scriptsize\sffamily}]
  \node[rectangle, rounded corners=3pt, draw=FRMGold!70, fill=PaleGold, line width=0.8pt,
    inner sep=6pt, align=center] (c) at (0,0) {\textbf{Value of}\\ \textbf{convexity} $\uparrow$};
  \node[rectangle, rounded corners=3pt, draw=FRMNavy!55, fill=PaleNavy, line width=0.8pt,
    inner sep=6pt] (v) at (5.0,0.85) {\textbf{Volatility} $\uparrow$};
  \node[rectangle, rounded corners=3pt, draw=FRMNavy!55, fill=PaleNavy, line width=0.8pt,
    inner sep=6pt] (m) at (5.0,-0.85) {\textbf{Maturity} $\uparrow$};
  \draw[FRMGold, -{Stealth[length=4.5pt]}, line width=0.9pt] (c) -- (v);
  \draw[FRMGold, -{Stealth[length=4.5pt]}, line width=0.9pt] (c) -- (m);
  \node[anchor=south, text=black!60] at (2.6,0.42) {is affected by};
\end{tikzpicture}
\end{center}
\end{mrkeybox}

% ---------------------------------------------------------
\lo{MR-13 d}{Identify the components into which the return on a bond can be decomposed, and calculate the expected return on a bond for a risk-averse investor.}

\noindent\gapbadge\hspace{4pt}{\sffamily\fontsize{7.6}{9}\selectfont\color{black!55}%
The source notes cover the risk-premium calculation but never identify the three components of
bond return. That half of the objective is written below from the GARP curriculum.}\par
\vspace{6pt}

\begin{mrgapbox}[title={The three components of bond return}]
The instantaneous return on a zero-coupon bond decomposes into exactly three pieces:
\[ \frac{dP}{P} \;=\; \underbrace{f(T)\,dt}_{\text{passage of time}}
   \;-\; \underbrace{D\,dr}_{\text{rate change}}
   \;+\; \underbrace{\tfrac{1}{2}\,C\,\sigma^{2}\,dt}_{\text{convexity}} \]
\begin{itemize}
\item \term{Passage of time.} With the term structure unchanged, the bond earns the forward rate
      $f(T)$ corresponding to its term.
\item \term{Duration.} Increases in the rate reduce the bond's return in proportion to its
      duration $D$.
\item \term{Convexity.} The volatility of rates --- movement either up \textit{or} down ---
      increases return in proportion to convexity $C$. Across portfolios with the same duration,
      more convex portfolios gain more as rates move, whichever way they move.
\end{itemize}
Taking expectations replaces $dr$ with $E[dr]$ and leaves the other two terms unchanged.
\end{mrgapbox}

\subsection{Expected return for a risk-averse investor}

\begin{mrgapbox}[title={Risk-neutral versus risk-averse}]
\term{Risk-neutral} investors require no risk premium, so they demand that every bond offer an
expected return equal to the short-term rate. \term{Risk-averse} investors demand higher
expected returns for bonds with greater interest rate risk, and that risk is measured by
duration. The premium is proportional to duration and does \textit{not} depend on the
characteristics of any individual bond.
\end{mrgapbox}

\begin{mrfmlbox}
\[ E\!\left[\frac{dP}{P}\right] \;=\; r_0\,dt \;+\; \lambda\,D\,dt \]
\mrvk{$r_0$ = the short-term rate;\quad $\lambda$ = the risk premium per year of duration
risk;\quad $D$ = the bond's duration. The premium scales with duration, so a longer bond earns
more of it.}{}
\end{mrfmlbox}

\begin{mrexbox}[title={Example --- expected return with a duration-based premium}]
The short-term rate is 1\%, the duration of a bond is 5, and the risk premium is 10 basis points
per year of duration risk. What is the bond's expected return?
\tcblower
\small
\[ 1\% + \bigl(5 \times 0.1\%\bigr) \;=\; \mathbf{1.5\%} \text{ per year} \]
\end{mrexbox}

\subsection{Price and return with a risk premium, from the source notes}

\begin{mrdefbox}[title={Risk premium}]
The risk premium is the price for risk-averse investors to bear uncertainty.
\end{mrdefbox}

\begin{mrexbox}[title={Calculate the price of the 2-year zero with a 20 basis point risk premium}]
The rate tree is 8\% today, moving to 10\% or 6\% with equal probability. Add 20 basis points to
each one-year rate at the next node.
\tcblower
\small
\[ P \;=\; \frac{\left[ \dfrac{\$1}{1.102} + \dfrac{\$1}{1.062} \right] \times 0.5}{1.08}
   \;=\; \frac{\bigl[\$0.90744 + \$0.94162\bigr] \times 0.5}{1.08}
   \;=\; \mathbf{\$0.85605} \]
\end{mrexbox}

\begin{mrexbox}[title={Calculate the return on the 2-year zero with a 20 basis point risk premium}]
\small
The bond is bought today at \$0.85605. In one year it will be worth \$0.90909 or \$0.94340 with
equal probability, using the \textit{un}-premium-adjusted rates.
\[ \text{expected return} \;=\;
   \frac{\bigl[\$0.90909 + \$0.94340\bigr] \times 0.5 - \$0.85605}{\$0.85605}
   \;=\; \mathbf{0.082 \text{, or } 8.2\%} \]
\end{mrexbox}

\begin{mrkeybox}[title={Read the result}]
The one-year rate is 8\%, and the expected return on the two-year bond is 8.2\%. The extra 20
basis points is exactly the risk premium, which is what a risk-averse investor demanded for
bearing the extra period of rate uncertainty.
\end{mrkeybox}

\begin{center}
\begin{tikzpicture}[every node/.style={font=\scriptsize\sffamily}]
  \fill[PaleNavy, draw=FRMNavy!60, line width=0.7pt] (0,0) rectangle (9.6,0.8);
  \fill[PaleGold, draw=FRMGold!70, line width=0.7pt] (9.6,0) rectangle (10.2,0.8);
  \node[text=FRMNavy] at (4.8,0.4) {\textbf{short rate 8.00\%}};
  \node[anchor=west, text=FRMGold] at (10.4,0.4) {\textbf{$+$ 0.20\% risk premium}};
  \node[anchor=north, text=black!60] at (5.1,-0.3)
    {total expected return on the 2-year zero $=$ \textbf{8.20\%}};
\end{tikzpicture}
\end{center}
\figcap{The risk premium is additive on top of the short rate. Nothing else changes.}


% ==========================================================
